% Generated draft: requires mathematical and source review.
\documentclass[11pt]{article}
\usepackage[margin=25mm]{geometry}
\usepackage{fontspec}
\setmainfont{Latin Modern Roman}
\usepackage{amsmath,amssymb,mathtools,unicode-math}
\setmathfont{Latin Modern Math}
\usepackage{enumitem,needspace,xurl,hyperref}
\setlength{\parindent}{0pt}
\setlength{\parskip}{5pt}
\setlength{\emergencystretch}{3em}
\newcommand{\N}{\mathbb{N}}
\newcommand{\Z}{\mathbb{Z}}
\newcommand{\Q}{\mathbb{Q}}
\newcommand{\R}{\mathbb{R}}
\newcommand{\C}{\mathbb{C}}
\newcommand{\F}{\mathbb{F}}
\newcommand{\A}{\mathbb{A}}
\newcommand{\PP}{\mathbb{P}}
\newcommand{\E}{\mathbb{E}}
\newcommand{\eps}{\varepsilon}
\newcommand{\dd}{\,\mathrm{d}}
\begin{document}
\section*{UnsolvedMath ID: 30002336}
\textbf{Source problem code:} catalogue entry
\subsection{Definitions and mathematical statement}
The supplied record collects several distinct questions. The following conventional definitions clarify the ropelength questions; the unresolved conventions and missing definitions are stated explicitly.

Let a knot be an embedded closed curve \(\gamma\colon S^1\hookrightarrow\R^3\), and let \([\gamma]\) denote its ambient isotopy class. For a compact set \(A\subset\R^3\), define its reach by
\[
\operatorname{reach}(A)=\sup\{r>0:\text{ every }x\in\R^3\text{ with }\operatorname{dist}(x,A)<r\text{ has a unique nearest point in }A\}.
\]
For a rectifiable knot with positive reach, use the thickness convention \(\operatorname{Thi}(\gamma)=\operatorname{reach}(\gamma(S^1))\), and define
\[
\operatorname{Rop}(\gamma)=\frac{\operatorname{Len}(\gamma)}{\operatorname{Thi}(\gamma)},\qquad
\mathcal R(K)=\inf_{[\gamma]=K}\operatorname{Rop}(\gamma),
\]
where the infimum is over rectifiable positive-reach representatives. A representative attaining this infimum is ropelength-minimizing in \(K\). This is a conventional framework for expressing the questions; the record does not specify the admissible regularity class or thickness normalization. It also does not define ropelength-criticality or its admissible variations, so that term remains unresolved.

\begin{enumerate}
\item \textbf{Explicit representatives and trefoil regularity.} Can a ropelength-minimizing or ropelength-critical knot be described explicitly? For a ropelength-minimizing representative of the trefoil knot type, is it piecewise \(C^2\), or piecewise smooth? Here piecewise smooth ordinarily means a finite decomposition into smooth arcs, but the source does not specify endpoint regularity or whether smooth means \(C^\infty\). It does not specify whether the trefoil question concerns the existence of one such minimizer or every minimizer, and it does not identify a unique minimizing representative. The explicit-description question also leaves unspecified whether the knot must be nontrivial, what constitutes an explicit description, and which trefoil chirality is intended. The supplied literature assessment specifically refers to explicit nontrivial minimizers.

\item \textbf{Growth in crossing number.} The original statement additionally asks whether there are knot families whose minimum ropelength grows more than linearly in crossing number, and whether there is a better lower bound for alternating knots. Let \(c(K)\) denote the minimum crossing number over diagrams of \(K\); an alternating knot type is one admitting a diagram in which overcrossings and undercrossings alternate along the knot. One possible precise interpretation of the growth question is whether there exists a sequence \((K_j)_{j\geq1}\) such that
\[
c(K_j)\longrightarrow\infty,
\qquad
\frac{\mathcal R(K_j)}{c(K_j)}\longrightarrow\infty.
\]
This interpretation is provisional: the record does not prescribe a growth rate or specify whether crossing number means the invariant \(c(K)\) or the crossing count of selected diagrams. The alternating-knot question does not identify the lower bound to be improved or the intended improvement. These questions occur in the original statement but are omitted from the supplied clean statement; they are retained here without resolving that discrepancy.

\item \textbf{Higher-dimensional integral Menger energies.} Which geometric quantities or topological invariants of \(\Sigma\) can be bounded in terms of \(M_p(\Sigma)\)? The record supplies neither a formula for this energy nor the dimension and ambient space of \(\Sigma\), admissible regularity, range of \(p\), normalization, or auxiliary hypotheses. Thus \(M_p\) is retained as source notation, not identified with a particular integral Menger energy. The intended quantities and the permitted dependence of a bound on dimension, scale, size, or other parameters are also unspecified. No unique quantified energy inequality can be extracted from this record.

\item \textbf{Loop number and writhe.} For a regular oriented unknot diagram \(D\), let
\[
w(D)=\sum_{x\in\operatorname{Cross}(D)}\operatorname{sgn}(x)
\]
be its writhe, using the usual oriented crossing signs. If \(L(D)\) denotes the loop number intended by the source, is
\[
L(D)\geq |w(D)|\qquad\text{for every unknot diagram }D?
\]
What further relations connect loop number with known knot invariants? The source does not define loop number, specify allowed diagram conventions, or state whether any minimization over diagrams is involved. Consequently \(L(D)\) is an unresolved placeholder, and the displayed inequality is the supplied target rather than a fully defined conjecture. The additional requested relations and their quantifiers are not specified.
\end{enumerate}

These are statement-only questions. The supplied status and verification labels do not resolve the formulation gaps above or constitute independent mathematical review.
\subsection{Short English statement}
The record asks for explicit ropelength minimizers or critical knots, regularity of a minimizing trefoil, growth of ropelength relative to crossing number, bounds from higher-dimensional Menger energies, and relations involving loop number. In particular, it asks whether every unknot diagram satisfies \(L(D)\geq |w(D)|\). Definitions of criticality, the energy, and loop number are missing, and several quantifiers remain unclear. The supplied classification is partially solved; its dated assessment requires independent review.
\subsection{Sources}
\begin{itemize}
\item Local UnsolvedMath registry: supplied complete record for permanent ID \(30002336\), title ``Ropelength, Menger Energy, and Loop-Number Problems,'' source problem code \texttt{catalogue entry}, in the Oberwolfach Reports Open Problems collection. The proposed year is supplied as 2013; no individual proposer is named.
\item Supplied source citation: \emph{Geometric Knot Theory (2014)}, Oberwolfach Reports, DOI \texttt{10.4171/owr/2013/22}, pp.~1313--1358. \url{https://doi.org/10.4171/owr/2013/22}. The literature-source list also labels this reference ``Oberwolfach Report 22/2013.'' The supplied 2014 title date and 2013 report identification are preserved without reconciliation.
\item Cantarella--Kusner--Sullivan, \emph{On the Minimum Ropelength of Knots and Links}. Supplied literature link: \url{https://arxiv.org/abs/math/0103224}.
\item \emph{The ropelength conjecture of alternating knots}. Supplied literature link: \url{https://www.cambridge.org/core/journals/mathematical-proceedings-of-the-cambridge-philosophical-society/article/ropelength-conjecture-of-alternating-knots/3B154F4E32EC26108D2C38EDFD4B8AA6}.
\item Supplied literature assessment, checked \textbf{2026-08-21}: the record labels the collection of questions ``partially solved'' and states that minimizers exist, that several ropelength bounds and energy estimates are known, including a linear lower bound for alternating knots, and that explicit nontrivial minimizers, smoothness of the trefoil, and the remaining invariant questions are unresolved. This is supplied dated curation derived from the OWR collection, not an independently verified literature assessment.
\item The registry labels its statement \texttt{corrected\_verified} and reports a check against the official Oberwolfach report. That verification claim is supplied provenance only; the sources have not been independently consulted for this draft.
\end{itemize}
\end{document}
